\documentclass[11pt]{article}
\usepackage[margin=1in]{geometry}
\usepackage{amsmath,amssymb,amsthm}
\usepackage[hidelinks]{hyperref}
\newtheorem{theorem}{Theorem}
\title{An exact evaluation of the two-shifted q-Motzkin determinant}
\author{Principia Math In Progress}
\date{18 September 2026}
\begin{document}
\maketitle
\begin{abstract}
We evaluate the determinant for every size, using determinant condensation
and Cigler's established zero- and one-shift evaluations. This determines
the missing power of \(q\) and the complete normalized polynomial.
The printed conjecture omits a constant term when the size is divisible
by three. Campaign eligibility remains pending: the complete final
journal source has not been inspected. No new research credit is claimed.
\end{abstract}

\section*{Task 1: the original problem and its complete evaluation}
Let \(M_0(q)=1\) and
\[
 M_{n+1}(q)=M_n(q)+
       \sum_{r=0}^{n-1}q^{r+1}M_r(q)M_{n-r-1}(q).
\]
Put
\[
 D_n^{(s)}=\det[M_{i+j+s}(q)]_{i,j=0}^{n-1},
 \qquad D_0^{(s)}=1,\qquad
 S_n=\sum_{j=0}^{n-1}j^2=\frac{n(n-1)(2n-1)}6 .
\]
MathDB 339076 reproduces Conjecture A.2 on p.~62 of
\cite{OS}. It predicts \(D_n^{(2)}=q^{c_n}f_n(q)\) with
\[
 f_n(q)=
 \begin{cases}
 \displaystyle\sum_{\substack{1\le l\le n\\l\not\equiv1\pmod3}}q^l,
                    &n\equiv0\pmod3,\\[5pt]
 \displaystyle(1+q)\sum_{r=0}^{\lfloor n/3\rfloor}q^{3r},
                    &n\not\equiv0\pmod3.
 \end{cases}
\]
The lower bound \(1\le l\) in the first branch is also present in the
primary preprint. The following theorem evaluates the whole family,
not merely the first exception to that printed formula.

\begin{theorem}\label{evaluation}
For every integer \(n\ge0\), the following polynomial identity holds:
\begin{equation}\label{sum}
 D_n^{(2)}
   =q^{S_n+n^2}
      \sum_{\substack{0\le j\le n\\j\not\equiv2\pmod3}}q^{-j}.
\end{equation}
For \(n\ge1\), equivalently \(D_n^{(2)}=q^{c_n}g_n(q)\), where
\begin{align}
 c_n&=S_n+n^2-n+
             \begin{cases}1,&n\equiv2\pmod3,\\0,&\text{otherwise},\end{cases}
                                                        \label{exponent}\\
 g_n(q)&=
 \begin{cases}
 \displaystyle\sum_{\substack{0\le l\le n\\l\not\equiv1\pmod3}}q^l,
                    &n\equiv0\pmod3,\\[5pt]
 \displaystyle(1+q)\sum_{r=0}^{\lfloor n/3\rfloor}q^{3r},
                    &n\not\equiv0\pmod3.
 \end{cases}                                            \label{normalized}
\end{align}
In particular, every nonzero coefficient of \(g_n\) is \(1\);
\(g_n(0)=1\); and \(c_n\) is the exact order of vanishing at \(q=0\).
\end{theorem}

\subsection*{A condensation proof}
Cigler's established formulas \cite[Eqs.~(39)--(40)]{Cigler} are
\begin{equation}\label{known}
 D_n^{(0)}=q^{S_n},\qquad
 D_n^{(1)}=\kappa_n q^{S_n+n(n-1)/2},
 \qquad
 (\kappa_n)_{n\ge0}=(1,1,0,-1,-1,0,1,1,0,\ldots).
\end{equation}
The period is six. These two evaluations are credited prior inputs,
not new assertions of this note. We work first in the field
\(\mathbb Q(q)\), where each \(D_n^{(0)}\) is nonzero.

The Desnanot--Jacobi identity, applied to the \((n+1)\)-square
unshifted Hankel matrix, gives for every \(n\ge1\)
\begin{equation}\label{condense}
 D_{n+1}^{(0)}D_{n-1}^{(2)}
      =D_n^{(0)}D_n^{(2)}-(D_n^{(1)})^2 .
\end{equation}
Indeed, the two principal corner minors have shifts zero and two;
each off-diagonal corner minor has shift one; and deleting both
extreme rows and columns gives the size \(n-1\), shift-two minor.
This also applies to \(n=1\), with the empty determinant equal to one.

Set \(E_n=D_n^{(2)}/D_{n+1}^{(0)}\). Dividing
\eqref{condense} by \(D_n^{(0)}D_{n+1}^{(0)}\) and using
\(S_{n+1}-S_n=n^2\), we obtain
\[
 E_n-E_{n-1}
   =\frac{(D_n^{(1)})^2}{D_n^{(0)}D_{n+1}^{(0)}}
   =\kappa_n^2q^{\,n(n-1)-n^2}
   =\kappa_n^2q^{-n}.
\]
Now \(E_0=1=\kappa_0^2\), and
\(\kappa_j^2=1\) precisely when \(j\not\equiv2\pmod3\).
Summing the recurrence yields
\[
 E_n=\sum_{\substack{0\le j\le n\\j\not\equiv2\pmod3}}q^{-j}.
\]
Multiplying by \(D_{n+1}^{(0)}=q^{S_n+n^2}\) proves \eqref{sum}.
No division by a one-shift determinant occurs; its periodic zeros
therefore cause no exception.

To normalize the result, the largest permitted index is \(n\) when
\(n\not\equiv2\pmod3\), and \(n-1\) otherwise.
Reversing the summation about this largest index gives
\eqref{exponent}--\eqref{normalized}. The exponent \(c_n\) is nonnegative
because \(S_n\ge0\) and \(n^2-n\ge0\). Both sides of the evaluated
identity are polynomials, so the equality obtained in
\(\mathbb Q(q)\) is also valid at \(q=0\) and at every complex \(q\).
This proves the theorem. \qed

\subsection*{The printed formula, scope, and checks}
For example,
\[
 D_1^{(2)}=1+q,\qquad D_2^{(2)}=q^4(1+q),\qquad
 D_3^{(2)}=q^{11}(1+q^2+q^3).
\]
Thus the printed \(f_3=q^2+q^3\) cannot be repaired by changing only
its monomial prefactor: it has two nonzero terms, while \(D_3^{(2)}\)
has three. For every positive multiple of three one must add the
constant term to \(f_n\). The theorem supplies the full corrected
evaluation for all sizes; a transcription correction alone would not
qualify as a new campaign solution.

The accompanying verifier builds the moments directly from their
defining recurrence and computes exact polynomial determinants.
It checks all three shifts at sizes \(0\le n\le16\), including zeros
in the one-shift sequence, and compares the two forms of the answer.
These finite checks audit the indexing and exponents. The proof for
unbounded \(n\) is the condensation argument above.

The 2019 final article \cite{OSfinal} was identified bibliographically
but could not be obtained in full. The current preprint was inspected
selectively, especially Appendix A; it was not read throughout.
Consequently the source and priority audit is incomplete, and this
package has zero new campaign credit. An authoring-agent audit and
exact checks do not constitute independent peer review.

\section*{Task 2: a direct weighted-path interpretation}
Construct a weight-preserving description of the terms
\(q^{S_n+n^2-j}\), \(j\not\equiv2\pmod3\), directly from the
shifted Hankel path families. The objective is an explicit cancellation
map whose fixed points explain the missing residue class.

\section*{Task 3: the three-shift determinant}
Evaluate \(D_n^{(3)}\) for every size and determine its sign and
normalized factors in all three residue classes. Start from condensation
with shift one; handle the vanishing \(D_n^{(1)}\) without dividing by
zero. This is a proposed direction, not a solved additional problem.

\section*{Task 4: higher even shifts}
Prove or disprove coefficientwise positivity of \(D_n^{(2k)}\) for
all positive \(n,k\), the separate Conjecture A.4 in \cite{OS}.
The present theorem covers only \(k=1\). Positivity of a Gram
determinant for real \(q>0\) does not establish coefficientwise
positivity.

\section*{Task 5: retain a horizontal-step parameter}
Replace the moment recursion by
\[
 M_{n+1}(a,q)=aM_n(a,q)+
       \sum_{r=0}^{n-1}q^{r+1}M_r(a,q)M_{n-r-1}(a,q).
\]
Find the corresponding shift-two determinant and identify values of
\(a\) giving periodic normalized one-shift determinants. Determine
which specializations retain nonnegative coefficients in \(q\).

\begin{thebibliography}{9}
\bibitem{Cigler}
J. Cigler, \emph{q-Catalan- und q-Motzkinzahlen},
Sitzungsber.\ Abt.\ II \textbf{208} (1999), 3--20,
especially Eqs.~(37), (39), (40), printed pp.~15--16.
\url{https://homepage.univie.ac.at/johann.cigler/preprints/cat-motzkin.pdf}.
\bibitem{OS}
S.-j. Oh and T. Scrimshaw, \emph{Identities from representation theory},
arXiv:1805.00113v1 (2018), Appendix A, pp.~61--63.
\url{https://arxiv.org/abs/1805.00113v1}.
\bibitem{OSfinal}
S.-j. Oh and T. Scrimshaw, \emph{Identities from representation theory},
Discrete Math.\ \textbf{342}(9) (2019), 2493--2541.
\url{https://doi.org/10.1016/j.disc.2019.05.020}.
Bibliographic identification only; complete final text not inspected.
\end{thebibliography}
\end{document}
